% spring_math.tex -- equations implemented by springcontroller.
%
% A fragment, not a standalone document: \input{} it from the paper.
% Requires: amsmath, amssymb, bm
%
% Every equation here corresponds to code in
% springcontroller/springcontroller/virtual_spring.py; the marginal
% comments give the class and method.

\section{Virtual spring control}
\label{sec:spring-math}

The arm is controlled by a set of virtual springs acting in parallel. Each
spring is a passive element that maps the arm's current configuration to a
generalised joint torque; the controller sums their contributions and commands
the result in torque mode. No inverse kinematics and no trajectory are
involved, which is what allows a person to push the arm off its setpoint and
feel a smooth restoring force rather than a servo fighting back.

\subsection{Notation and kinematics}
\label{sec:notation}

Let $\bm{q} \in \mathbb{R}^{n}$ be the vector of joint positions and
$\dot{\bm{q}} \in \mathbb{R}^{n}$ the joint velocities, for an arm with $n$
degrees of freedom. Forward kinematics gives, for each link $\ell$, a
homogeneous transform from the link frame to the world frame,
%
\begin{equation}
  \bm{T}_{w\ell}(\bm{q}) =
  \begin{bmatrix}
    \bm{R}_{\ell}(\bm{q}) & \bm{o}_{\ell}(\bm{q}) \\
    \bm{0}^{\top} & 1
  \end{bmatrix},
  \qquad
  \bm{R}_{\ell} \in SO(3), \quad \bm{o}_{\ell} \in \mathbb{R}^{3}.
\end{equation}

A spring attaches to link $\ell$ at a point ${}^{\ell}\bm{a} \in \mathbb{R}^{3}$
fixed in the \emph{link's own} frame. Its world position is
%
\begin{equation}
  \bm{p}(\bm{q}) \;=\; \bm{o}_{\ell}(\bm{q}) \;+\; \bm{R}_{\ell}(\bm{q})\,
  {}^{\ell}\bm{a}.
  \label{eq:attachment-point}
\end{equation}
%
Specifying the attachment in the link frame rather than the world frame is what
lets a spring track a point on a moving link---for instance the centre of a
block held in the gripper, which is offset from any URDF frame origin.

Throughout, $\hat{\bm{v}} = \bm{v}/\lVert\bm{v}\rVert$ denotes a unit vector,
$[\bm{v}]_{\times}$ the skew-symmetric matrix satisfying
$[\bm{v}]_{\times}\bm{w} = \bm{v}\times\bm{w}$, and $\bm{A}^{+}$ the
Moore--Penrose pseudoinverse of $\bm{A}$. Displacement vectors are written to
point \emph{from} the current state \emph{toward} the target, so that a spring's
force carries a positive stiffness and is attractive.

\subsection{Jacobian and the transpose mapping}
\label{sec:jacobian}

The geometric Jacobian of the attachment point,
$\bm{J}(\bm{q}) \in \mathbb{R}^{6 \times n}$, relates joint velocity to the
point's linear velocity and the link's angular velocity:
%
\begin{equation}
  \begin{bmatrix} \dot{\bm{p}} \\ \bm{\omega} \end{bmatrix}
  \;=\;
  \bm{J}(\bm{q})\,\dot{\bm{q}},
  \qquad
  \bm{J} = \begin{bmatrix} \bm{J}_{v} \\ \bm{J}_{\omega} \end{bmatrix},
  \qquad
  \bm{J}_{v}, \bm{J}_{\omega} \in \mathbb{R}^{3 \times n}.
  \label{eq:jacobian-def}
\end{equation}

Kinematics libraries report the Jacobian at a link's frame \emph{origin}. Since
the attachment point is generally offset from that origin, the translational
block must be shifted. Writing the offset in world-aligned axes as
$\bm{r} = \bm{R}_{\ell}\,{}^{\ell}\bm{a}$, the attachment point is rigidly
fixed to the link, so its velocity is that of the origin plus the rotational
contribution,
%
\begin{equation}
  \dot{\bm{p}} \;=\; \dot{\bm{o}}_{\ell} + \bm{\omega} \times \bm{r}
  \;=\; \dot{\bm{o}}_{\ell} - [\bm{r}]_{\times}\,\bm{\omega},
\end{equation}
%
and therefore, with $\bm{J}_{v}^{o}$ the translational Jacobian at the origin,
%
\begin{equation}
  \boxed{\;\bm{J}_{v}(\bm{q}) \;=\; \bm{J}_{v}^{o}(\bm{q}) \;-\;
  [\bm{r}]_{\times}\,\bm{J}_{\omega}(\bm{q}).\;}
  \label{eq:point-shift}
\end{equation}
%
The rotational block is unchanged by the shift, since every point of a rigid
link shares the same angular velocity.

Springs are defined as wrenches---a force $\bm{f}$ applied at the attachment
point and a moment $\bm{m}$ on the link---but the arm is commanded in joint
torques. The two are related by the principle of virtual work. An arbitrary
virtual joint displacement $\delta\bm{q}$ produces, by
\eqref{eq:jacobian-def}, a virtual translation
$\delta\bm{p} = \bm{J}_{v}\,\delta\bm{q}$ and a virtual rotation
$\delta\bm{\phi} = \bm{J}_{\omega}\,\delta\bm{q}$. Equating the work done in
Cartesian space with the work done in joint space,
%
\begin{equation}
  \delta W
  \;=\; \bm{f}^{\top}\delta\bm{p} + \bm{m}^{\top}\delta\bm{\phi}
  \;=\; \bm{f}^{\top}\bm{J}_{v}\,\delta\bm{q}
      + \bm{m}^{\top}\bm{J}_{\omega}\,\delta\bm{q}
  \;=\; \bigl(\bm{J}_{v}^{\top}\bm{f}
      + \bm{J}_{\omega}^{\top}\bm{m}\bigr)^{\!\top}\delta\bm{q}
  \;\stackrel{!}{=}\; \bm{\tau}^{\top}\delta\bm{q}.
\end{equation}
%
As this must hold for every $\delta\bm{q}$, the bracketed vectors are equal:
%
\begin{equation}
  \boxed{\;\bm{\tau} \;=\; \bm{J}^{\top}
  \begin{bmatrix}\bm{f}\\\bm{m}\end{bmatrix}
  \;=\; \bm{J}_{v}^{\top}\bm{f} \;+\; \bm{J}_{\omega}^{\top}\bm{m}.\;}
  \label{eq:jacobian-transpose}
\end{equation}

Equation~\eqref{eq:jacobian-transpose} is the whole of the mapping from spring
to actuator, and it has two properties worth stating because the rest of the
design depends on them. First, it requires no matrix inversion and so is
defined everywhere, including at kinematic singularities---unlike an
inverse-kinematics or resolved-rate scheme, it cannot diverge as the arm
straightens. Second, it is exactly the static force balance: the commanded
torques are those that would hold the arm in equilibrium against the spring
wrench, which is why the resulting behaviour is passive and feels compliant
rather than driven.

Column $j$ of $\bm{J}_{v}$ is the linear velocity the attachment point acquires
per unit velocity of joint $j$, so the $j$th entry of
$\bm{J}_{v}^{\top}\bm{f}$ is the projection of the spring force onto that
direction. A joint that cannot move the attachment point therefore receives no
torque from a force spring, no matter how stiff---a fact returned to in
\S\ref{sec:joint-spring}.

\subsection{Linear springs}
\label{sec:linear-spring}

A linear spring pulls the attachment point $\bm{p}(\bm{q})$ toward a fixed
world target $\bm{p}^{*}$. Define the displacement, its magnitude (the
\emph{extension}), and its direction:
%
\begin{equation}
  \bm{d}(\bm{q}) = \bm{p}^{*} - \bm{p}(\bm{q}),
  \qquad
  e = \lVert\bm{d}\rVert,
  \qquad
  \hat{\bm{d}} = \bm{d}/e \quad (e > 0).
\end{equation}

Rather than a single Hooke's-law term, the elastic force is piecewise in $e$,
so that a tolerance region can be specified in which the arm is left entirely
free. Given a stiffness $k > 0$, an inner radius $r_{\mathrm{in}}$, an outer
radius $r_{\mathrm{out}}$, and a rest length $L_{0}$, and writing $\beta$ for
the predicate $r_{\mathrm{out}} > r_{\mathrm{in}}$ (whether a deadband is
configured at all):
%
\begin{equation}
  \bm{f}_{k}(e) =
  \begin{cases}
    \bm{0}
      & e \le r_{\mathrm{in}}, \\[6pt]
    \dfrac{k\,(e - r_{\mathrm{in}})^{2}}{r_{\mathrm{out}} - r_{\mathrm{in}}}\,
    \hat{\bm{d}}
      & \beta \wedge \; r_{\mathrm{in}} < e \le r_{\mathrm{out}}, \\[10pt]
    k\,(e - r_{\mathrm{in}})\,\hat{\bm{d}}
      & \beta \wedge \; e > r_{\mathrm{out}}, \\[6pt]
    k\,(e - L_{0})\,\hat{\bm{d}}
      & \neg\beta .
  \end{cases}
  \label{eq:linear-force}
\end{equation}

The three-region structure gives a dead zone, a soft-start ramp, and a
fully-stiff region. Inside $r_{\mathrm{in}}$ the spring is silent: the operator
may move the arm freely within a tolerance ball around the target. Across the
deadband the force grows \emph{quadratically}, entering from zero rather than
stepping to full stiffness at the boundary. Beyond $r_{\mathrm{out}}$ the
spring is an ordinary linear spring anchored at $r_{\mathrm{in}}$.

Anchoring the outer region at $r_{\mathrm{in}}$ rather than at $L_{0}$ is
deliberate: it makes \eqref{eq:linear-force} continuous at
$e = r_{\mathrm{out}}$ by construction, both the ramp and the linear branch
evaluating to $k\,(r_{\mathrm{out}} - r_{\mathrm{in}})$ there, for any
configured $L_{0}$. Were the outer branch anchored at $L_{0}$ instead, a
mismatched rest length would put a force discontinuity exactly at
$r_{\mathrm{out}}$, and sensor noise carrying $e$ back and forth across that
boundary would make the commanded force chatter between the two values on
successive control cycles. Note that \eqref{eq:linear-force} is $C^{0}$ but not
$C^{1}$: the slope at $e = r_{\mathrm{out}}$ jumps from $2k$ on the ramp to $k$
on the linear branch. This is a discontinuity in stiffness, not in force, and
is not perceptible in practice.

Viscous damping opposes the world-space velocity of the attachment point, which
\eqref{eq:jacobian-def} supplies directly from the measured joint velocities
without numerical differentiation of position:
%
\begin{equation}
  \bm{f}_{b} \;=\; -\,b\,\dot{\bm{p}} \;=\; -\,b\,\bm{J}_{v}\dot{\bm{q}},
  \qquad b \ge 0.
  \label{eq:linear-damping}
\end{equation}

Applying \eqref{eq:jacobian-transpose} with zero moment, the spring's
contribution to the joint torques is
%
\begin{equation}
  \boxed{\;\bm{\tau} \;=\; \bm{J}_{v}^{\top}\bigl(\bm{f}_{k} + \bm{f}_{b}\bigr).\;}
  \label{eq:linear-torque}
\end{equation}

In the ungated case ($\neg\beta$) the elastic term is conservative, with stored
potential energy
%
\begin{equation}
  U = \tfrac{1}{2}\,k\,(e - L_{0})^{2}.
\end{equation}

\subsection{Pose springs}
\label{sec:pose-spring}

A pose spring orients a link rather than positioning it. A direction
${}^{\ell}\hat{\bm{n}}$ fixed in the link frame---the ``face normal'', for
example the outward normal of the face of a held block---is driven to point at
a world target $\bm{p}^{*}$, giving a \emph{look-at} constraint. In the study
this keeps a held object's face turned toward the participant as the arm moves.

The current face normal in world coordinates is
%
\begin{equation}
  \bm{n}(\bm{q}) \;=\; \bm{R}_{\ell}(\bm{q})\,{}^{\ell}\hat{\bm{n}},
\end{equation}
%
and the desired heading is the unit vector from the attachment point to the
target,
%
\begin{equation}
  \bm{u}(\bm{q}) \;=\;
  \frac{\bm{p}^{*} - \bm{p}(\bm{q})}{\lVert\bm{p}^{*} - \bm{p}(\bm{q})\rVert}.
  \label{eq:pose-desired}
\end{equation}
%
Both are recomputed from the \emph{current} configuration on every control
cycle. This is what distinguishes a look-at constraint from a fixed orientation
setpoint: as the arm moves, $\bm{u}$ moves with it, and the face continues to
track $\bm{p}^{*}$ rather than holding a rotation frozen at the moment the
spring was created.

The rotation carrying $\bm{n}$ onto $\bm{u}$ is expressed in axis--angle form
via their cross product:
%
\begin{equation}
  \bm{c} = \bm{n} \times \bm{u},
  \qquad
  \theta = \operatorname{atan2}\bigl(\lVert\bm{c}\rVert,\;
                                     \bm{n}\cdot\bm{u}\bigr),
  \qquad
  \hat{\bm{a}} = \bm{c}/\lVert\bm{c}\rVert.
  \label{eq:pose-axis-angle}
\end{equation}
%
The two-argument form is used in preference to
$\theta = \arccos(\bm{n}\cdot\bm{u})$ because it remains well conditioned at
both ends of the range: $\arccos$ has unbounded derivative as its argument
approaches $\pm 1$, so near perfect alignment ($\theta \to 0$) and near
inversion ($\theta \to \pi$)---precisely the configurations the controller
spends most of its time in---it amplifies floating-point error in the dot
product into a large angular error. The axis is undefined when
$\lVert\bm{c}\rVert = 0$; the implementation takes $\hat{\bm{a}} = \bm{0}$
there, yielding zero moment. At $\theta = 0$ this is correct. At $\theta = \pi$
it is an \emph{unstable} equilibrium---the exactly-antiparallel configuration
is left unforced, but any perturbation off it produces a moment driving the
normal back toward alignment, so the arm does not rest there in practice.

The restoring moment is Hooke's law in the angle, with viscous damping opposing
the link's angular velocity:
%
\begin{equation}
  \bm{m} \;=\; \underbrace{k_{\theta}\,\theta\,\hat{\bm{a}}}_{\text{elastic}}
  \;-\; \underbrace{b_{\theta}\,\bm{J}_{\omega}\dot{\bm{q}}}_{\text{damping}},
  \qquad k_{\theta}, b_{\theta} \ge 0,
  \label{eq:pose-moment}
\end{equation}
%
which by \eqref{eq:jacobian-transpose}, with zero force, gives the
orientation torque
%
\begin{equation}
  \bm{\tau}_{\mathrm{orient}} \;=\; \bm{J}_{\omega}^{\top}\bm{m}.
  \label{eq:pose-orient-torque}
\end{equation}

\subsubsection{Bounding positional drift}
\label{sec:pose-nullspace}

Equation~\eqref{eq:pose-orient-torque} carries no Cartesian force, but it does
not follow that the attachment point stays put. On a serial chain nearly every
joint affects position and orientation simultaneously, so torque applied to
correct orientation will in general also translate the point---potentially far
outside the workspace a paired linear spring is trying to hold. The pose spring
therefore splits its torque according to whether each component can move the
attachment point.

Let $\bm{N}$ be the orthogonal projector onto the null space of $\bm{J}_{v}$,
%
\begin{equation}
  \bm{N} \;=\; \bm{I}_{n} - \bm{J}_{v}^{+}\bm{J}_{v},
\end{equation}
%
and decompose accordingly:
%
\begin{equation}
  \bm{\tau}_{\mathrm{safe}} = \bm{N}\,\bm{\tau}_{\mathrm{orient}},
  \qquad
  \bm{\tau}_{\mathrm{risky}} = \bm{\tau}_{\mathrm{orient}}
                             - \bm{\tau}_{\mathrm{safe}}
                             = \bm{J}_{v}^{+}\bm{J}_{v}\,
                               \bm{\tau}_{\mathrm{orient}}.
  \label{eq:pose-split}
\end{equation}
%
By the Moore--Penrose identity $\bm{J}_{v}\bm{J}_{v}^{+}\bm{J}_{v} =
\bm{J}_{v}$ we have $\bm{J}_{v}\bm{\tau}_{\mathrm{safe}} = \bm{0}$ exactly:
the safe component provably cannot translate the attachment point, and is
applied at full strength unconditionally. Since $\bm{N}$ is symmetric and
idempotent, the two components are orthogonal and the split is a genuine
decomposition rather than an approximation.

The remaining component is attenuated by how far the attachment point has
drifted from a nominal centre $\bm{c}$, within a permitted radius $\rho > 0$.
With $\delta = \lVert\bm{p}(\bm{q}) - \bm{c}\rVert$,
%
\begin{equation}
  \alpha = \min\!\left(1,\; \frac{\rho}{\delta}\right),
  \qquad
  \bm{\tau}_{\mathrm{pose}} = \bm{\tau}_{\mathrm{safe}}
                            + \alpha\,\bm{\tau}_{\mathrm{risky}}.
  \label{eq:pose-gate}
\end{equation}
%
Inside the radius $\alpha = 1$ and the spring has full authority; beyond it
authority decays as $\rho/\delta$. The gate is continuous at $\delta = \rho$,
so no chattering arises as the point crosses the boundary, and it needs no
separate margin parameter. Discarding $\bm{\tau}_{\mathrm{risky}}$ entirely
was found to be too weak in practice: a real arm's null space is small or
empty in many configurations, leaving almost no orientation authority at all.

Because \eqref{eq:pose-gate} only ever \emph{throttles} torque and never
generates a restoring pull, a pose spring used without a paired linear spring
has nothing returning the attachment point once it drifts. An optional
translational term supplies this, zero inside $\rho$ and Hooke's law beyond it:
%
\begin{equation}
  \bm{f}_{p} =
  \begin{cases}
    \bm{0} & \delta \le \rho, \\[4pt]
    k_{p}\,(\delta - \rho)\,\dfrac{\bm{c} - \bm{p}(\bm{q})}{\delta}
           & \delta > \rho,
  \end{cases}
  \qquad k_{p} \ge 0,
  \label{eq:pose-position-force}
\end{equation}
%
which shares the ``excess distance past a radius'' form of the outer branch of
\eqref{eq:linear-force}. The complete pose spring torque is then
%
\begin{equation}
  \boxed{\;\bm{\tau} \;=\; \bm{\tau}_{\mathrm{safe}}
                        + \alpha\,\bm{\tau}_{\mathrm{risky}}
                        + \bm{J}_{v}^{\top}\bm{f}_{p}.\;}
  \label{eq:pose-torque}
\end{equation}

\subsection{Joint springs}
\label{sec:joint-spring}

The third spring type acts on a single joint's own coordinate, with no
Jacobian involved. For joint $j$ with angle $\theta_{j}$ and target
$\theta^{*}$,
%
\begin{equation}
  \boxed{\;\tau_{j} \;=\; -\,k\,\Delta\theta \;-\; b\,\dot{\theta}_{j},
  \qquad
  \Delta\theta = \operatorname{wrap}_{(-\pi,\pi]}\!\bigl(\theta_{j}
                 - \theta^{*}\bigr),\;}
  \label{eq:joint-torque}
\end{equation}
%
contributing $\tau_{j}\bm{e}_{j}$ to the joint torque vector, where
$\bm{e}_{j}$ is the $j$th standard basis vector. The sign is explicit here
because $\Delta\theta$ is measured as current minus target, the opposite
convention to $\bm{d}$ in \S\ref{sec:linear-spring}. Taking the shortest signed
path avoids driving a continuous joint the long way round when the error
straddles the wrap point.

This element exists because of the structural limitation noted at the end of
\S\ref{sec:jacobian}. Consider a wrist joint whose axis passes through the
attachment point. Rotating it does not move that point at all, so the
corresponding column of $\bm{J}_{v}$ is zero, and by
\eqref{eq:linear-torque} a linear spring produces \emph{exactly} zero torque on
that joint regardless of its stiffness---leaving the joint free to drift. A
joint spring supplies restoring or damping torque on precisely those axes a
Cartesian spring is blind to.

\subsection{Aggregation}
\label{sec:aggregation}

Springs superpose. For a set of $N$ active springs with individual torques
$\bm{\tau}_{i}$ from \eqref{eq:linear-torque}, \eqref{eq:pose-torque} or
\eqref{eq:joint-torque}, the commanded torque is
%
\begin{equation}
  \boxed{\;\bm{\tau}_{\mathrm{cmd}} \;=\;
  s\sum_{i=1}^{N}\bm{\tau}_{i}(\bm{q}, \dot{\bm{q}})
  \;+\; \bm{\tau}_{g}(\bm{q}),\;}
  \qquad s \in [0, 1],
  \label{eq:total-torque}
\end{equation}
%
where $\bm{\tau}_{g}(\bm{q})$ is the generalised gravity torque and $s$ is a
ramp factor used to fade springs in when torque control is engaged, avoiding a
step in commanded torque at the transition.

The scale $s$ multiplies the spring sum only. Gravity compensation is applied
at full magnitude on every cycle in which it is enabled, since it is what holds
the arm up; ramping it alongside the springs would let the arm sag as $s$
approaches zero. On platforms that compensate gravity in firmware,
$\bm{\tau}_{g} = \bm{0}$ and the term drops out.
